\subsection{Обратный инжиниринг CAD-моделей: современное состояние}

\subsubsection{Определение и классификация задач обратного инжиниринга}

Обратный инжиниринг (Reverse Engineering) в контексте CAD — это процесс создания цифровой модели изделия по его физическому образцу или по имеющемуся цифровому представлению без исходной параметрической информации \cite{tulmash2025}. В более широком смысле, обратный инжиниринг включает восстановление не только геометрии, но и конструкторской документации, а также понимание принципов работы и технологии изготовления изделия \cite{interpolyaris2022}.

В машиностроении задачи обратного инжиниринга возникают в следующих ситуациях \cite{ritm2025}:
\begin{itemize}
\item замена утраченных или изношенных деталей, когда закупка запчастей невозможна;
\item воссоздание изделий, снятых с производства;
\item импортозамещение — создание отечественных аналогов импортных комплектующих;
\item модернизация существующих изделий при отсутствии исходной документации;
\item анализ и улучшение конструкции на основе физического образца.
\end{itemize}

Традиционные методы обратного инжиниринга включают 3D-сканирование физического объекта с последующей обработкой облака точек и созданием CAD-модели. Этот процесс трудоёмок и требует высокой квалификации оператора. Однако в данной работе рассматривается частный, но не менее важный случай: обратный инжиниринг по уже имеющемуся цифровому представлению (STEP-файлу) — задача восстановления параметрической истории без физического образца.

\subsubsection{Традиционные алгоритмические методы и их ограничения}

Классические подходы к восстановлению истории построения основаны на анализе B-Rep представления и поиске геометрических примитивов. Например, для обнаружения операции выдавливания алгоритм может искать набор параллельных граней одинаковой формы, соединённых боковыми поверхностями. Для обнаружения скругления — искать рёбра с цилиндрическими или тороидальными сопряжениями.

Однако такие методы имеют ряд фундаментальных ограничений:

Неоднозначность интерпретации. Одна и та же итоговая геометрия может быть получена разными последовательностями операций. Например, ступенчатый вал может быть создан несколькими последовательными выдавливаниями или одной операцией вращения с последующими вырезаниями. Алгоритмический метод, как правило, находит только одно решение, которое может не соответствовать исходному замыслу конструктора.

Чувствительность к порядку операций. При изменении порядка операций итоговая геометрия может остаться той же, но промежуточные состояния будут различаться. Традиционные методы обычно не учитывают эту вариативность.

Сложность обработки «нечистых» моделей. Импортированные модели часто содержат артефакты: «битые» грани, микро-зазоры, невалидную топологию. Это затрудняет работу алгоритмов, основанных на точной топологии.

Отсутствие обучаемости. Алгоритмические методы не улучшаются с накоплением опыта — каждую новую модель они анализируют «с нуля», не используя знания о ранее встреченных паттернах.

\subsubsection{Подход Autodesk Research и других лабораторий}

В 2020 году Autodesk Research опубликовала результаты первых экспериментов по применению машинного обучения для восстановления параметрической истории \cite{autodesk2020}. Исследователи использовали датасет Fusion 360 Gallery — коллекцию из тысяч моделей, созданных пользователями в среде Fusion 360, с сохранением полной истории построения.

В рамках узкой постановки задачи — восстановление последовательности операций «эскиз + выдавливание» для простых моделей — были получены обнадёживающие результаты. Нейросетевая модель обучалась предсказывать, какие грани нужно выдавить на каждом шаге. Однако, как отмечают авторы, «даже лучшие модели машинного обучения ошибаются примерно раз на сто попыток. На более сложных моделях мы всё ещё не можем восстановить исходную модель за разумное время» \cite{autodesk2020}.

Для преодоления этой проблемы исследователи применили гибридный подход: нейросеть используется для ранжирования гипотез (какие грани наиболее вероятно были выдавливанием), а затем выполняется детерминированный поиск по наиболее вероятным вариантам. Такой подход сочетает обучаемость нейросети с гарантией корректности детерминированных алгоритмов.

В более поздних работах исследователи расширили область применения, включив другие типы операций и альтернативные представления геометрии (сетки и облака точек). Однако, как подчёркивается в публикации, «обратный инжиниринг CAD-моделей — это захватывающая область для исследований в области машинного обучения, и мы только начинаем» \cite{autodesk2020}.

\subsubsection{Промышленные решения и их недостатки}

На рынке представлено несколько коммерческих продуктов для обратного инжиниринга, наиболее известные из которых — Geomagic Design X (3D Systems) и РеВерсия (ТЕСИС) \cite{reversia2024}. Эти продукты в основном ориентированы на обработку данных 3D-сканирования (облака точек и полигональные сетки) и позволяют создавать параметрические CAD-модели по физическим образцам.

Однако существующие промышленные решения имеют ряд ограничений применительно к задаче восстановления истории по STEP-файлу:
\begin{itemize}
\item они требуют значительного участия человека (интерактивная сегментация, назначение типов поверхностей);
\item результаты сильно зависят от квалификации оператора;
\item они не «понимают» логику исходного конструктора — восстанавливаемая история может быть неоптимальной или неестественной;
\item стоимость лицензий на профессиональное ПО для обратного инжиниринга высока.
\end{itemize}

Таким образом, разработка автоматизированного решения на основе машинного обучения, интегрированного непосредственно в среду КОМПАС-3D, является актуальной и востребованной задачей.

